\documentclass[11pt]{article}
\usepackage[a4paper,margin=27mm]{geometry}
\usepackage[T1]{fontenc}
\usepackage[utf8]{inputenc}
\usepackage{lmodern}
\usepackage{microtype}
\usepackage{amsmath,amssymb,amsthm,mathtools}
\usepackage{booktabs,longtable,array,tabularx}
\usepackage{enumitem}
\usepackage{xcolor}
\usepackage{hyperref}
\usepackage{listings}
\usepackage{fancyvrb}
\usepackage{enumitem}
\usepackage{url}
\hypersetup{colorlinks=true,linkcolor=blue!55!black,citecolor=blue!55!black,urlcolor=blue!55!black,pdftitle={A Constructive Atlas of Trigonometric Identities in Intuitionistic Primitive Recursive Arithmetic},pdfauthor={ProveIt research artifact},pdfsubject={Constructive trigonometric identities in iPRA}}
\setlist{nosep}
\newtheorem{theorem}{Theorem}[section]
\newtheorem{lemma}[theorem]{Lemma}
\newtheorem{proposition}[theorem]{Proposition}
\newtheorem{corollary}[theorem]{Corollary}
\theoremstyle{definition}
\newtheorem{definition}[theorem]{Definition}
\newtheorem{construction}[theorem]{Construction}
\newtheorem{remark}[theorem]{Remark}
\newcommand{\N}{\mathbb N}
\newcommand{\Z}{\mathbb Z}
\newcommand{\Q}{\mathbb Q}
\newcommand{\C}{\mathbb C}
\newcommand{\R}{\mathbb R}
\newcommand{\iPRA}{\mathsf{iPRA}}
\newcommand{\cis}{\operatorname{cis}}
\newcommand{\apart}{\mathrel{\#}}
\newcommand{\eps}{\varepsilon}
\newcommand{\angles}{\mathcal A}
\newcommand{\laur}{\mathcal L}
\newcommand{\den}{\operatorname{den}}
\newcommand{\num}{\operatorname{num}}
\newcommand{\atanTwo}{\operatorname{atan2}}
\newcommand{\abs}[1]{\left|#1\right|}
\newcommand{\norm}[1]{\left\lVert#1\right\rVert}
\newcommand{\floor}[1]{\left\lfloor#1\right\rfloor}
\newcommand{\ceil}[1]{\left\lceil#1\right\rceil}
\newcommand{\code}[1]{\texttt{\detokenize{#1}}}
\newcommand{\sem}[1]{[\![#1]\!]}
\DeclareMathOperator{\sgn}{sgn}
\DeclareMathOperator{\RePart}{Re}
\DeclareMathOperator{\ImPart}{Im}
\lstset{basicstyle=\ttfamily\small,breaklines=true,columns=fullflexible,frame=single}
\sloppy

\title{A Constructive Atlas of Trigonometric Identities in Intuitionistic Primitive Recursive Arithmetic}
\author{Research artifact prepared for the ProveIt project}
\date{15 November 2026}

\begin{document}
\maketitle

\begin{abstract}
This report gives a constructive proof architecture for the large majority of the
finite identities catalogued in Wikipedia's \emph{List of trigonometric identities},
and for several of its analytic and infinite families, inside intuitionistic
primitive recursive arithmetic (\(\iPRA\)).  Because ordinary \(\iPRA\) has only
natural-number objects, the phrase ``a theorem about real sine and cosine in
\(\iPRA\)'' is made precise by a definitional interpretation: a real is a fixed
primitive-recursive rapidly convergent dyadic name together with an arithmetic
regularity certificate.  Equality is extensional equality of names; strict order
and division use positive separation, or \emph{apartness}, witnesses.

The analytic core defines the complex exponential by factorial series with an
explicit primitive-recursive tail modulus.  The finite Cauchy-product identity
proves \(\exp(z+w)=\exp z\exp w\), hence the sine and cosine addition laws and the
Pythagorean identity.  A quarter turn is then constructed, without an appeal to a
classical intermediate-value theorem, by a bounded rational sign-search for the
unique zero of cosine in \([1,2]\); \(\pi\) is twice that zero.

Once this foundation is present, almost all elementary identities are finite
algebra.  We isolate a Gaussian Laurent-polynomial certificate calculus in which
\(\cis(x+y)=\cis x\,\cis y\) and \(\cis(\pi/2)=i\).  Its executable reference
implementation checks 227 identities in 27 families, with exact rational
arithmetic and zero failures.  Polynomial identities are unconditional.
Tangent, cotangent, secant, cosecant, square-root, phase, and inverse-function
forms are stated with the constructive side conditions that their usual printed
versions suppress.  The archive contains this report, the exact checker and its
certificate ledger, a coverage table, and a Lean implementation plan tied to the
ProveIt \code{ConciseIPRA} development.
\end{abstract}

\tableofcontents

\section{Scope, status, and the precise theorem}

The source list is not one theorem.  It combines polynomial identities, rational
identities with hidden domain restrictions, statements about signs and branches,
finite sums and products, power-series formulas, inverse functions, and several
substantial theorems about infinite products or Bessel functions.  A responsible
constructive treatment must separate these layers.

\begin{definition}[Three proof levels]
We use the following status terminology.
\begin{enumerate}[label=\textup{(\roman*)}]
\item A \emph{division-free identity} is an equality built from integer or
rational constants, addition, multiplication, sine, cosine, and integer angle
combinations.  It is asserted without side conditions.
\item A \emph{guarded rational identity} is asserted after every denominator is
supplied with an apartness witness.  Its cross-multiplied polynomial form is
usually unconditional.
\item An \emph{analytic identity} includes a convergence modulus, an order or
branch certificate, or both.  The modulus is required to be primitive recursive.
\end{enumerate}
\end{definition}

\begin{theorem}[Constructive trigonometric atlas]
\label{thm:atlas}
For every fixed finite tuple of primitive-recursive regular real names, \(\iPRA\)
proves the translations of all division-free identities in Sections
\ref{sec:algebra-families}--\ref{sec:triangles}.  For every fixed tuple of such
names and fixed apartness or branch certificates, \(\iPRA\) proves the guarded
translations in those sections.  It also proves the exponential, sine, cosine,
quarter-turn, inverse-function, differentiation, and dyadic-product theorems of
Sections \ref{sec:exp}--\ref{sec:pi}, \ref{sec:inverse}, and
\ref{sec:infinite}, with the explicit moduli given there.

The proof is uniform syntactically: from the codes of the displayed
primitive-recursive terms and their certificates one constructs a finite
\(\iPRA\) derivation by a primitive-recursive proof transformer.
\end{theorem}

The last sentence is deliberately not phrased as quantification over all
primitive-recursive functions.  There is no primitive-recursive universal
function evaluating every primitive-recursive program, and ordinary first-order
\(\iPRA\) has no function variables.  The theorem is therefore a proof schema
over fixed primitive-recursive terms, or equivalently a conservative definitional
extension by those terms.  This is exactly the form in which numerical analysis
is normally represented in a first-order arithmetic weak enough to care about
program extraction.

The archive's \code{artifacts/coverage.csv} gives a row-by-row status.  The main
boundary is as follows.

\begin{center}
\begin{tabularx}{\textwidth}{@{}>{\bfseries}p{0.24\textwidth}X@{}}
\toprule
Completed in this report & Basic, reciprocal and Pythagorean identities;
parity, reflections, shifts and periodicity; all sum/difference identities;
finite many-angle, multiple-angle and Chebyshev schemas; half-angle and
power-reduction formulas; product-to-sum and sum-to-product schemas; finite
geometric, Lagrange, Dirichlet and arithmetic-progression sums; roots-of-unity
products; finite Viete products; Hermite's cotangent identity; amplitude and
linear-fractional identities; Euler and hyperbolic conversions; principal
inverse compositions; triangle identities; derivatives of sine and cosine; and
the dyadic infinite cosine product.\\
\addlinespace
Completed with explicit guards & Every quotient form involving tangent,
cotangent, secant or cosecant; signed square-root half-angle formulas; phase
representations; inverse addition formulas; and formulas with poles in finite
sums or products.\\
\addlinespace
Constructive proof outline supplied & Jacobi--Anger/Bessel generating identities
and the broader termwise differentiation/integration library.\\
\addlinespace
Not claimed complete & The full Euler--Weierstrass product over all zeros of
sine and application-specific medical-imaging operator identities.  These need
substantial product or integration libraries beyond the trigonometric core.\\
\bottomrule
\end{tabularx}
\end{center}

\section{What ``in \(\iPRA\)'' means}
\label{sec:meaning}

\subsection{The arithmetic base}

The ProveIt \code{ConciseIPRA} development presents the usual first-order
intuitionistic primitive recursive arithmetic with a symbol for every primitive
recursive function, defining equations, successor axioms, and induction for
quantifier-free formulas.  The Lean file \code{ConciseIPRA/PR.lean} represents a
primitive-recursive symbol of arity \(k\) by an inductive type \code{PR k}; its
constructors are zero, successor, projection, composition, and primitive
recursion.  The arithmetic axioms contain the two recursion equations and
quantifier-free induction.  We use exactly that interface.

No real numbers are added as primitive objects.  All real-number language below
is an abbreviation for formulas about natural-number codes.

\subsection{Signed integers, dyadics, and rapid names}

A signed integer is represented by a pair \((a,b)\in\N^2\), denoting \(a-b\).
A dyadic is a triple \((a,b,e)\), denoting \((a-b)/2^e\).  Addition,
multiplication, comparison after a common denominator, absolute value, and
rounding to a prescribed denominator are primitive recursive.

\begin{definition}[Rapid dyadic name]
A rapid name is a fixed primitive-recursive dyadic sequence
\(p\mapsto x[p]\), supplied with an \(\iPRA\) proof of
\[
  \abs{x[p+1]-x[p]}\le 2^{-p-3}.
  \tag{2.1}
\]
It also carries an integer magnitude bound \(B\ge1\) such that
\(\abs{x[p]}\le B\) for every \(p\).
\end{definition}

The geometric sum proves, intuitionistically and with only finite induction,
that for \(q\ge p\),
\[
 \abs{x[q]-x[p]}\le \sum_{j=p}^{q-1}2^{-j-3}<2^{-p-2}.
 \tag{2.2}
\]
Thus the closed interval
\[
 I_x(p)=[x[p]-2^{-p-2},\ x[p]+2^{-p-2}]
 \tag{2.3}
\]
contains every later approximation.  One may regularize the intervals by taking
finite intersections, but none of the proofs needs that cosmetic step.

\begin{definition}[Extensional equality and apartness]
For rapid names \(x,y\), write \(x=_\R y\) when
\[
  \forall p\;\abs{x[p+4]-y[p+4]}\le 2^{-p}.
  \tag{2.4}
\]
An apartness witness \(x\apart0\) is a finite tuple
\((p,m,\sigma)\), \(\sigma\in\{-1,+1\}\), certifying that the whole interval
\(I_x(p)\) lies respectively above \(2^{-m}\) or below \(-2^{-m}\).
Likewise, \(x<y\) is used only with a finite precision at which
\(I_x(p)\) and \(I_y(p)\) are strictly separated.
\end{definition}

Apartness is positive information.  From \(x\ne0\) in intuitionistic logic one
cannot in general compute a reciprocal; from \(x\apart0\) one can.  This is the
reason every quotient identity later has an explicit guard.

\subsection{Arithmetic on names}

Negation and addition are pointwise with a fixed precision shift.  To multiply
\(x\) and \(y\), use bounds \(B_x,B_y\), request the input intervals at a
precision exceeding the output precision by
\(\ceil{\log_2(B_x+B_y+2)}+O(1)\), multiply their endpoints, and round outward.
All loops are bounded by data already present in the names.  The regularity
proof is the elementary estimate
\[
 \abs{xy-x'y'}\le \abs{x}\abs{y-y'}+\abs{y'}\abs{x-x'}.
 \tag{2.5}
\]

If an apartness witness gives \(0<d\le\abs{x}\), interval inversion is applied
at a precision for which the input width is below \(d/2\).  The estimate
\[
 \abs{x^{-1}-y^{-1}}=\frac{\abs{x-y}}{\abs{x}\abs{y}}
 \le \frac{2\abs{x-y}}{d^2}
 \tag{2.6}
\]
gives a primitive-recursive precision shift.  Square root on a nonnegative
bounded interval is obtained by a fixed number of rational bisections; if a
positive lower bound is available it is Lipschitz, and at zero the elementary
modulus \(\abs{\sqrt u-\sqrt v}\le\sqrt{\abs{u-v}}\) is used.  The latter merely
doubles the requested input precision.

\begin{proposition}[Closure]
Rapid names with certificates are closed under negation, addition, bounded
multiplication, reciprocal from apartness, nonnegative square root, finite sums,
and finite products.  Codes for all these operations are primitive recursive,
and their regularity proofs use only quantifier-free induction.
\end{proposition}

\begin{proof}
Every construction is a bounded loop on dyadic integers.  Equations (2.2),
(2.5), and (2.6) provide quantifier-free loop invariants.  Universal closure in
the precision parameter is introduced after proving the corresponding fixed
precision inequality.  No choice sequence or undecidable sign test is used.
\end{proof}

\subsection{The translation of a real statement}

A displayed equality \(F(x_1,\ldots,x_r)=G(x_1,\ldots,x_r)\) denotes the
first-order formula obtained by replacing every operation with its name
constructor and replacing equality by (2.4).  A theorem ``for all constructive
reals'' means: for arbitrary fixed \code{PR} terms for the input names and
arbitrary fixed derivations of their regularity assumptions, a uniform proof
transformer emits an \(\iPRA\) proof of the translated equality.  This avoids
both second-order quantification and any hidden classical completion of \(\Q\).

\section{The complex exponential with an explicit primitive-recursive modulus}
\label{sec:exp}

We use the \(\ell^1\) norm on Gaussian rationals,
\(\norm{a+ib}_1=\abs a+\abs b\), because it is rational and submultiplicative.
Complex rapid names are pairs of real names.

For a complex number bounded by an integer \(B\ge1\), define the finite sum
\[
 E_N(z)=\sum_{n=0}^{N-1}\frac{z^n}{n!}.
 \tag{3.1}
\]
Factorials, powers, and bounded sums are primitive recursive.

\begin{lemma}[Explicit exponential tail]
\label{lem:exp-tail}
Let \(\norm z_1\le B\), \(B\ge1\), and put
\[
 N(B,p)=2B^2+2B+p+4.
 \tag{3.2}
\]
Then
\[
 \sum_{n=N(B,p)}^\infty\frac{\norm z_1^n}{n!}<2^{-p-3}.
 \tag{3.3}
\]
The proof of the corresponding finite-tail bound is formalizable in \(\iPRA\).
\end{lemma}

\begin{proof}
For \(n\ge2B\), successive majorant terms have ratio
\(B/(n+1)\le1/2\), so the tail is at most twice its first term.  Write
\(N=2B+m\), where \(m=2B^2+p+4\).  Discarding the positive factor
\((2B)!\) and using \(B\le2^B\),
\[
 2\frac{B^N}{N!}
 \le 2B^{2B}2^{-m}
 \le 2\,2^{2B^2}2^{-(2B^2+p+4)}=2^{-p-3}.
\]
All inequalities concern natural-number products and powers.  The ratio-tail
argument is a finite induction followed by a geometric-sum estimate, so
quantifier-free induction suffices.
\end{proof}

To evaluate \(E_N\) at a real name, request a sufficiently fine input interval
and perform outward interval arithmetic.  A deliberately coarse sensitivity
bound is obtained by replacing every coefficient and every intermediate value
by \((B+1)^N\); requesting
\[
 p+(B+2)N(B,p)+N(B,p)+8
 \tag{3.4}
\]
input bits is more than sufficient.  The exact shift is unimportant; its
primitive recursiveness and proved adequacy are what matter.

\begin{definition}[Constructive exponential, sine, and cosine]
The name \(\exp z\) at precision \(p\) is the outward-rounded value of
\(E_{N(B,p+4)}(z)\).  Put
\[
 \sin x=\frac{\exp(ix)-\exp(-ix)}{2i},\qquad
 \cos x=\frac{\exp(ix)+\exp(-ix)}2,
 \qquad \cis x=\exp(ix).
 \tag{3.5}
\]
Division by \(2\) is exact dyadic division and needs no apartness guard.
\end{definition}

\subsection{The multiplication law}

The key finite identity is
\[
 \sum_{r+s=k}\frac{x^r}{r!}\frac{y^s}{s!}
 =\frac1{k!}\sum_{r=0}^k\binom{k}{r}x^r y^{k-r}
 =\frac{(x+y)^k}{k!}.
 \tag{3.6}
\]
It is a polynomial identity over \(\Z\) after multiplication by \(k!\).
The binomial theorem and the factorial identity
\(\binom{k}{r}r!(k-r)!=k!\) are proved by primitive recursion.

\begin{theorem}[Exponential addition]
\label{thm:exp-add}
For bounded complex rapid names,
\[
 \exp(x+y)=\exp x\,\exp y.
 \tag{3.7}
\]
\end{theorem}

\begin{proof}
Multiply finite sums and group terms of total degree \(k\).  Equation (3.6)
identifies every coefficient below the truncation boundary.  The discarded
parts are bounded by Lemma \ref{lem:exp-tail}; the product error is bounded by
the same tails multiplied by the elementary bound
\(\norm{\exp z}_1\le 2^{2B+2}\).  Given an output precision \(p\), all
truncation indices and input precisions are supplied by explicit primitive-
recursive expressions.  Hence \(\iPRA\) proves the precision-\(p\) equality;
universal introduction gives (3.7).
\end{proof}

Taking \(y=-x\) gives \(\exp x\exp(-x)=1\).  Substituting \(ix,iy\) and taking
real and imaginary parts gives the fundamental formulas
\begin{align}
 \sin(x+y)&=\sin x\cos y+\cos x\sin y,\tag{3.8}\\
 \cos(x+y)&=\cos x\cos y-\sin x\sin y.\tag{3.9}
\end{align}
Moreover,
\[
 \cos^2x+\sin^2x
 =\frac{(u+u^{-1})^2-(u-u^{-1})^2}{4}=1,
 \qquad u=\cis x.
 \tag{3.10}
\]
These are the only analytic facts needed by the finite algebraic atlas.

\section{A primitive-recursive construction of the quarter turn and \(\pi\)}
\label{sec:pi}

We cannot define \(\pi\) by saying ``choose the least positive zero of sine.''
That would hide both a choice principle and a real-number minimization.  Instead
we construct a rapidly convergent rational sequence with a verified error
bound.

\subsection{Alternating estimates}

For \(0\le t\le2\), the absolute terms in the sine and cosine series decrease
from the first nonconstant term onward.  The alternating-series estimates,
proved by finite pairing and the explicit tail bound, yield
\begin{align}
 t-\frac{t^3}{6}&\le\sin t\le t,\tag{4.1}\\
 \cos1&\ge1-\frac12=\frac12,\tag{4.2}\\
 \cos2&\le1-2+\frac{16}{24}=-\frac13.\tag{4.3}
\end{align}
In particular, for \(0\le t\le2\),
\[
 \sin t\ge t\left(1-\frac{t^2}{6}\right)\ge\frac t3.
 \tag{4.4}
\]

From the already proved addition formulas,
\[
 \cos a-\cos b
 =2\sin\frac{a+b}{2}\sin\frac{b-a}{2}.
 \tag{4.5}
\]
Thus for \(1\le a<b\le2\), equation (4.4) gives the quantitative strict
monotonicity estimate
\[
 \cos a-\cos b\ge\frac{b-a}{9}.
 \tag{4.6}
\]
The upper estimate \(\abs{\cos a-\cos b}\le\abs{a-b}\) follows from
\(\abs{\sin t}\le\abs t\).  Both bounds are algebraic consequences of finite
series estimates and require no derivative theorem.

\subsection{Located sign-search}

Fix an output precision \(p\).  Begin with the rational bracket \([1,2]\), whose
endpoint cosine signs are certified by (4.2)--(4.3).  Repeatedly take the
rational midpoint \(m\) and compute an interval for \(\cos m\) of width at most
\(2^{-p-12}\).  Rational comparison gives three decidable cases.
\begin{enumerate}[label=\textup{(\alph*)}]
\item If the interval lies strictly above zero, replace the left endpoint by
\(m\).
\item If it lies strictly below zero, replace the right endpoint by \(m\).
\item If it contains zero, record \(m\) as an approximate zero.
\end{enumerate}
Run at most \(p+12\) midpoint steps.  If case (c) occurs, then
\(\abs{\cos m}\le2^{-p-12}\).  If it never occurs, the final bracket has width
at most \(2^{-p-12}\); monotonicity and the Lipschitz estimate imply that its
midpoint \(m\) satisfies the same bound, up to one harmless bit.

Let \(r[p]\) be the resulting midpoint.  Equation (4.6) implies
\[
 \abs{r[p]-r[q]}
 \le9\bigl(\abs{\cos r[p]}+\abs{\cos r[q]}\bigr),
 \tag{4.7}
\]
so a fixed precision shift turns \(r\) into a rapid name.  This is a bounded
loop using only rational comparisons, hence a primitive-recursive construction.

\begin{definition}[Quarter turn and \(\pi\)]
Let \(q\) be the real represented by the shifted name \(r\), and define
\(\pi=2q\).
\end{definition}

\begin{theorem}[Quarter-turn specification]
\label{thm:quarter}
The theory \(\iPRA\) proves
\[
 1\le q\le2,\qquad \cos q=0,\qquad \sin q=1.
 \tag{4.8}
\]
Consequently
\[
 \sin\pi=0,\quad \cos\pi=-1,\quad
 \sin(2\pi)=0,\quad \cos(2\pi)=1,
 \tag{4.9}
\]
and sine and cosine have period \(2\pi\).
\end{theorem}

\begin{proof}
The construction and the exponential continuity modulus show
\(\cos q=0\).  Since \(q\in[1,2]\), (4.4) gives \(\sin q\ge1/3\).  Equation
(3.10) gives \((\sin q)^2=1\).  Factor
\((\sin q-1)(\sin q+1)=0\); the second factor is positively apart from zero,
so constructive cancellation gives \(\sin q=1\).  Equations (4.9) and
periodicity are now finite applications of (3.8)--(3.9).
\end{proof}

This construction is also useful in formalization: it prevents every later
identity involving \(\pi\) from depending on a large pre-existing constructive
real-analysis library.

\section{The finite Gaussian Laurent certificate calculus}
\label{sec:certificate}

The analytic work is now complete.  The rest of most of the source list can be
verified in a finite ring.

\begin{definition}[Angle syntax]
For angle variables \(x_1,\ldots,x_r\), an angle code is
\[
 a=(m_1,\ldots,m_r;k)\in\Z^r\times\Z,
 \qquad
 \sem{a}=\sum_jm_jx_j+kq.
 \tag{5.1}
\]
Addition and negation are componentwise.
\end{definition}

Let
\[
 \laur=\Q(i)[z_1^{\pm1},\ldots,z_r^{\pm1}]
 \tag{5.2}
\]
be the Gaussian Laurent ring.  Define
\begin{align}
 U(a)&=i^k\prod_jz_j^{m_j},\tag{5.3}\\
 S(a)&=\frac{U(a)-U(-a)}{2i},\qquad
 C(a)=\frac{U(a)+U(-a)}2.\tag{5.4}
\end{align}
A Laurent polynomial is represented by a finite list of monomials with Gaussian
rational coefficients.  Sorting monomials, merging equal exponents, and deleting
zero coefficients is a primitive-recursive normalizer.

\begin{theorem}[Certificate soundness]
\label{thm:certificate-sound}
If the normalized Laurent residue of two division-free expressions is zero,
then \(\iPRA\) proves their equality under every assignment of the angle
variables to regular real names.

For rational expressions \(P/Q\) and \(R/T\), a zero normalized residue
\(PT-RQ\) proves equality whenever the evaluations of \(Q\) and \(T\) are
apart from zero.
\end{theorem}

\begin{proof}
Evaluate \(z_j\) as \(\cis x_j\).  Theorem \ref{thm:exp-add} and Theorem
\ref{thm:quarter} give
\(U(a+b)=U(a)U(b)\), \(U(-a)=U(a)^{-1}\), and \(U(q)=i\).  Equations (5.4)
then evaluate to the constructed sine and cosine.  Soundness of each finite
normalizer step is ring arithmetic; induction on the finite term and monomial
lists is quantifier-free induction.  The rational case is cross multiplication
followed by constructive cancellation using the supplied apartness witnesses.
\end{proof}

The accompanying Python implementation uses only exact fractions and finite
maps; it performs no floating-point trigonometry.  Its certificate ledger
contains 227 zero residues in 27 families.  The finite checks are regressions,
not a replacement for the schema proofs: for example, the report proves the
\(n\)-angle binomial theorem for arbitrary \(n\), while the executable ledger
checks \(0\le n\le10\).

\section{Basic, symmetry, periodicity, and addition families}
\label{sec:algebra-families}

Put \(s_x=\sin x\), \(c_x=\cos x\).  The reciprocal functions are partial
constructive functions:
\[
 \tan x=\frac{s_x}{c_x}\quad(c_x\apart0),\quad
 \cot x=\frac{c_x}{s_x}\quad(s_x\apart0),\quad
 \sec x=\frac1{c_x},\quad \csc x=\frac1{s_x}.
 \tag{6.1}
\]
Their familiar Pythagorean consequences are
\[
 1+\tan^2x=\sec^2x,
 \qquad 1+\cot^2x=\csc^2x,
 \tag{6.2}
\]
with the corresponding guards.  Cross multiplication reduces both to
\(s_x^2+c_x^2=1\).

\subsection{Parity derived constructively}

It is worth deriving parity rather than assuming it, because the derivation
illustrates constructive cancellation without division.  Setting \(y=-x\) in
the addition laws gives
\begin{align*}
 s_xc_{-x}+c_xs_{-x}&=0,\\
 c_xc_{-x}-s_xs_{-x}&=1.
\end{align*}
Multiply the first equation by \(c_x\), the second by \(s_x\), and subtract.
Using \(s_x^2+c_x^2=1\) gives \(s_{-x}=-s_x\).  The analogous linear
combination gives \(c_{-x}=c_x\).  No case split on whether \(s_x\) or \(c_x\)
is zero is used.

The quarter-turn values then give all reflection and shift formulas, for example
\begin{align}
 \sin(q-x)&=\cos x,& \cos(q-x)&=\sin x,\tag{6.3}\\
 \sin(x+q)&=\cos x,& \cos(x+q)&=-\sin x,\tag{6.4}\\
 \sin(\pi-x)&=\sin x,& \cos(\pi-x)&=-\cos x,\tag{6.5}\\
 \sin(x+\pi)&=-\sin x,& \cos(x+\pi)&=-\cos x,\tag{6.6}\\
 \sin(x+2\pi)&=\sin x,& \cos(x+2\pi)&=\cos x.\tag{6.7}
\end{align}
Sign tables on quadrants are constructive consequences of (4.4), the shift
formulas, and a supplied strict interval locating the angle modulo \(2\pi\).
At an axis one states the exact zero separately instead of invoking an
undecidable trichotomy for reals.

\subsection{Sums and differences}

The unconditional formulas are
\begin{align}
 \sin(x\pm y)&=\sin x\cos y\pm\cos x\sin y,\tag{6.8}\\
 \cos(x\pm y)&=\cos x\cos y\mp\sin x\sin y.\tag{6.9}
\end{align}
For example, the tangent formula is best recorded first in cross-multiplied
form
\[
 \cos(x+y)\bigl(\tan x+\tan y\bigr)
 =\sin(x+y)\bigl(1-\tan x\tan y\bigr),
 \tag{6.10}
\]
with \(\cos x\apart0\) and \(\cos y\apart0\).  If in addition
\(1-\tan x\tan y\apart0\), cancellation yields
\[
 \tan(x+y)=\frac{\tan x+\tan y}{1-\tan x\tan y}.
 \tag{6.11}
\]
The corresponding formulas for subtraction, cotangent, secant, and cosecant
follow identically.  The exact checker includes all six common quotient forms.
The constructive correction is not a change in algebra; it is the replacement
of an implicit ``where defined'' by finite positive data sufficient to perform
the divisions.

\section{Multiple angles, Chebyshev polynomials, and finite angle sums}
\label{sec:multiple}

From \(\cis x=c_x+is_x\) and \(\cis(nx)=\cis(x)^n\), the binomial theorem gives,
for every \(n\in\N\),
\begin{align}
 \cos(nx)
 &=\sum_{j=0}^{\floor{n/2}}(-1)^j\binom{n}{2j}
 c_x^{n-2j}s_x^{2j},\tag{7.1}\\
 \sin(nx)
 &=\sum_{j=0}^{\floor{(n-1)/2}}(-1)^j\binom{n}{2j+1}
 c_x^{n-2j-1}s_x^{2j+1}.\tag{7.2}
\end{align}
This proof is a bounded binomial expansion and is directly available in
\(\iPRA\).  The familiar special cases include
\begin{align*}
 \sin2x&=2s_xc_x,&
 \cos2x&=c_x^2-s_x^2=2c_x^2-1=1-2s_x^2,\\
 \sin3x&=3s_x-4s_x^3,&
 \cos3x&=4c_x^3-3c_x.
\end{align*}
Quotient versions for \(\tan,\cot,\sec,\csc\) are obtained by guarded
cancellation.  For instance
\[
 \tan3x=\frac{3t-t^3}{1-3t^2},\qquad t=\tan x,
 \tag{7.3}
\]
provided the displayed terms are defined and the denominator is apart from
zero.

Define the Chebyshev polynomials by primitive recursion:
\begin{align*}
 T_0(X)&=1,&T_1(X)&=X,&T_{n+1}(X)&=2XT_n(X)-T_{n-1}(X),\\
 U_0(X)&=1,&U_1(X)&=2X,&U_{n+1}(X)&=2XU_n(X)-U_{n-1}(X).
\end{align*}
Adding the formulas for \(\cos((n+1)x)\) and \(\cos((n-1)x)\) proves
\[
 \cos(nx)=T_n(\cos x),\qquad
 \sin(nx)=\sin x\,U_{n-1}(\cos x).
 \tag{7.4}
\]
The induction invariant is a polynomial equality and therefore quantifier-free.

\subsection{A finite sum of arbitrary angles}

Let \(t_j=\tan x_j\), and let \(e_k\) be the elementary symmetric polynomial
of degree \(k\) in \(t_1,\ldots,t_n\), with \(e_0=1\).  Expanding
\[
 \prod_{j=1}^n(1+it_j)
 =(e_0-e_2+e_4-\cdots)+i(e_1-e_3+e_5-\cdots)
 \tag{7.5}
\]
and comparing it with
\(\prod_j\sec x_j\,\cis(\sum_jx_j)\) gives
\begin{align}
 \tan\sum_jx_j&=
 \frac{e_1-e_3+e_5-\cdots}{e_0-e_2+e_4-\cdots},\tag{7.6}\\
 \sec\sum_jx_j&=
 \frac{\prod_j\sec x_j}{e_0-e_2+e_4-\cdots},\tag{7.7}\\
 \csc\sum_jx_j&=
 \frac{\prod_j\sec x_j}{e_1-e_3+e_5-\cdots}.
 \tag{7.8}
\end{align}
Each quotient carries the obvious apartness guard.  The expansion is a finite
induction updating the elementary symmetric list from high degree to low degree;
the exact checker validates the first four arities, while (7.5) is the general
proof.

\section{Half angles, tangent half-angle substitution, and power reduction}
\label{sec:half-power}

The double-angle identities immediately give the unconditional square formulas
\[
 \sin^2\frac x2=\frac{1-\cos x}{2},\qquad
 \cos^2\frac x2=\frac{1+\cos x}{2}.
 \tag{8.1}
\]
A printed formula with a square root must include branch information:
\[
 \sin\frac x2=\sigma_s\sqrt{\frac{1-\cos x}{2}},\qquad
 \cos\frac x2=\sigma_c\sqrt{\frac{1+\cos x}{2}},
 \tag{8.2}
\]
where \(\sigma_s,\sigma_c\in\{-1,+1\}\) are certified by locating \(x/2\) in a
strict half-period interval.  At a zero endpoint, the value is stated as zero.
This avoids deciding the sign of an arbitrary real.

With the relevant denominators apart from zero,
\begin{align}
 \tan\frac x2&=\frac{\sin x}{1+\cos x}
 =\frac{1-\cos x}{\sin x},\tag{8.3}\\
 \cot\frac x2&=\frac{\sin x}{1-\cos x}
 =\frac{1+\cos x}{\sin x}.\tag{8.4}
\end{align}
Putting \(t=\tan(x/2)\) yields the Weierstrass substitution
\begin{align}
 \sin x&=\frac{2t}{1+t^2},&
 \cos x&=\frac{1-t^2}{1+t^2},&
 \cis x&=\frac{1+it}{1-it}.
 \tag{8.5}
\end{align}
The denominator \(1+t^2\) has a constructive positive lower bound \(1\), so
only the definition of \(t\) itself needs a cosine apartness witness.

\subsection{General power reduction}

The two compact master identities are
\begin{align}
 \cos^n x
 &=2^{-n}\sum_{k=0}^n\binom nk\cis((n-2k)x)
 =2^{-n}\sum_{k=0}^n\binom nk\cos((n-2k)x),\tag{8.6}\\
 \sin^n x
 &=(2i)^{-n}\sum_{k=0}^n(-1)^k\binom nk\cis((n-2k)x).
 \tag{8.7}
\end{align}
They are simply the binomial expansions of
\(((u+u^{-1})/2)^n\) and \(((u-u^{-1})/(2i))^n\).  Pairing terms \(k\) and
\(n-k\) gives the usual purely real even/odd tables.  For example,
\begin{align*}
 \sin^2x&=\frac{1-\cos2x}{2},&
 \cos^2x&=\frac{1+\cos2x}{2},\\
 \sin^3x&=\frac{3\sin x-\sin3x}{4},&
 \cos^3x&=\frac{3\cos x+\cos3x}{4},\\
 \sin^4x&=\frac{3-4\cos2x+\cos4x}{8},&
 \cos^4x&=\frac{3+4\cos2x+\cos4x}{8}.
\end{align*}
All divisions are by positive integers and are total on dyadic/rational names.

\section{Products to sums, sums to products, and finite products}
\label{sec:products}

Adding and subtracting the angle formulas proves
\begin{align}
 \cos x\cos y&=\frac{\cos(x-y)+\cos(x+y)}2,\tag{9.1}\\
 \sin x\sin y&=\frac{\cos(x-y)-\cos(x+y)}2,\tag{9.2}\\
 \sin x\cos y&=\frac{\sin(x+y)+\sin(x-y)}2,\tag{9.3}\\
 \cos x\sin y&=\frac{\sin(x+y)-\sin(x-y)}2.\tag{9.4}
\end{align}
Conversely, setting \(p=(x+y)/2\), \(r=(x-y)/2\) gives all four
sum-to-product formulas.  Division by two is rational scalar multiplication,
not a partial real division.

For arbitrary \(n\), finite expansion gives
\begin{align}
 \prod_{j=1}^n\cos x_j
 &=2^{-n}\sum_{\eps\in\{-1,1\}^n}
   \cos\left(\sum_j\eps_jx_j\right),\tag{9.5}\\
 \prod_{j=1}^n\sin x_j
 &=(2i)^{-n}\sum_{\eps\in\{-1,1\}^n}
   \left(\prod_j\eps_j\right)
   \cis\left(\sum_j\eps_jx_j\right).
 \tag{9.6}
\end{align}
The sum ranges over a primitive-recursively enumerable finite cube.  Pairing a
sign vector with its negative converts (9.6) into the familiar real sine or
cosine formula according to the parity of \(n\).

\subsection{Roots of unity}

For \(n\ge2\), let \(\zeta=\cis(2\pi/n)\).  Finite polynomial division gives
\[
 \frac{Z^n-1}{Z-1}=\prod_{k=1}^{n-1}(Z-\zeta^k)
 =1+Z+\cdots+Z^{n-1}.
 \tag{9.7}
\]
This is a finite identity in a cyclotomic extension, which can be represented
by pairs of polynomials modulo \(1+Z+\cdots+Z^{n-1}\); no algebraic closure is
needed.  Setting \(Z=1\), pairing conjugate factors, and using positivity of
sine on \((0,\pi)\) yields
\[
 \prod_{k=1}^{n-1}\sin\frac{k\pi}{n}=\frac{n}{2^{n-1}}.
 \tag{9.8}
\]
Likewise,
\[
 \prod_{k=0}^{n-1}\sin\left(x+\frac{k\pi}{n}\right)
 =2^{1-n}\sin(nx).
 \tag{9.9}
\]
Equation (9.9) is first proved after squaring or as a complex Laurent identity;
the sign is fixed by a rational sample interval.  Half-products follow by
positive square root, for example
\[
 \prod_{k=1}^{n-1}\sin\frac{k\pi}{2n}
 =\frac{\sqrt n}{2^{n-1}}.
 \tag{9.10}
\]
Every root sign in (9.8)--(9.10) is justified by an explicit interval
\(0<k\pi/n<\pi\).

\subsection{Finite Viete product}

Repeatedly applying \(\sin2t=2\sin t\cos t\) gives, for every \(N\ge1\),
\[
 2^N\sin x\prod_{j=0}^{N-1}\cos(2^jx)=\sin(2^Nx).
 \tag{9.11}
\]
This cross-multiplied form is unconditional.  Dividing by \(2^N\sin x\)
requires \(\sin x\apart0\).  The certificate ledger checks \(N\le9\), while
finite induction proves the schema.

\subsection{Hermite's cotangent identity}

If every \(\sin(a_k-a_j)\) for \(j\ne k\), and every displayed pole factor,
is apart from zero, then
\[
 \prod_{j=1}^n\cot(z-a_j)
 =\cos\frac{n\pi}{2}
 +\sum_{k=1}^n\cot(z-a_k)
   \prod_{j\ne k}\cot(a_k-a_j).
 \tag{9.12}
\]
For a human proof put \(w=e^{2iz}\), \(\alpha_k=e^{2ia_k}\), and use
\[
 \cot(z-a_k)=i\frac{w+\alpha_k}{w-\alpha_k}.
 \tag{9.13}
\]
Both sides of (9.12) are rational functions of \(w\).  At each simple pole
\(w=\alpha_k\) their principal parts agree, since the residue of the product is
the residue of its \(k\)-th factor times
\(\prod_{j\ne k}\cot(a_k-a_j)\).  Their difference is therefore constant.
Comparing leading coefficients after multiplication by
\(\prod_k(w-\alpha_k)\) gives \(\cos(n\pi/2)\).  This is finite polynomial
arithmetic.  The checker independently normalizes the cases \(n=2,3\).

\section{Finite trigonometric sums}
\label{sec:finite-sums}

The finite geometric identity
\[
 (1-r)\sum_{k=0}^nr^k=1-r^{n+1}
 \tag{10.1}
\]
is proved by a one-line induction in \(\iPRA\).  Put \(r=\cis h\) and use
\(1-e^{ih}=-2ie^{ih/2}\sin(h/2)\).  The unconditional cross-multiplied formula
is
\[
 \sin\frac h2\sum_{k=0}^n\cis(a+kh)
 =\sin\frac{(n+1)h}{2}\,
  \cis\left(a+\frac{nh}{2}\right).
 \tag{10.2}
\]
If \(\sin(h/2)\apart0\), division and real/imaginary parts give
\begin{align}
 \sum_{k=0}^n\cos(a+kh)
 &=\frac{\sin((n+1)h/2)}{\sin(h/2)}
   \cos\left(a+\frac{nh}{2}\right),\tag{10.3}\\
 \sum_{k=0}^n\sin(a+kh)
 &=\frac{\sin((n+1)h/2)}{\sin(h/2)}
   \sin\left(a+\frac{nh}{2}\right).
 \tag{10.4}
\end{align}
When the denominator is zero, (10.2) remains valid and the sum is handled by
the corresponding polynomial limiting case, such as \((n+1)\cos a\).

Useful specializations include
\begin{align}
 \sum_{k=0}^n\sin(2kx)
 &=\frac{\cos x-\cos((2n+1)x)}{2\sin x},\tag{10.5}\\
 \sum_{k=1}^n\cos(2kx)
 &=\frac{-\sin x+\sin((2n+1)x)}{2\sin x},\tag{10.6}\\
 1+2\sum_{k=1}^n\cos(2kx)
 &=\frac{\sin((2n+1)x)}{\sin x},\tag{10.7}\\
 \sum_{k=1}^n\cos((2k-1)x)
 &=\frac{\sin(2nx)}{2\sin x}.
 \tag{10.8}
\end{align}
Equation (10.7) is the Dirichlet-kernel identity.  Every quotient is guarded by
\(\sin x\apart0\); the cross-multiplied forms are unconditional.  The exact
ledger checks the Lagrange and Dirichlet families through \(n=10\), and the
arithmetic-progression formulas through \(n=8\).

\section{Linear combinations and linear-fractional rotations}
\label{sec:linear}

The identity
\[
 R\cos(x-\phi)=R\cos\phi\cos x+R\sin\phi\sin x
 \tag{11.1}
\]
is unconditional.  Hence, whenever data \(R,\phi\) satisfy
\(R\cos\phi=A\) and \(R\sin\phi=B\),
\[
 A\cos x+B\sin x=R\cos(x-\phi).
 \tag{11.2}
\]
One also has the unconditional amplitude equation
\[
 (A\cos u+B\cos v)^2+(A\sin u+B\sin v)^2
 =A^2+B^2+2AB\cos(u-v).
 \tag{11.3}
\]
For the usual canonical phase, let \(R=\sqrt{A^2+B^2}\).  If \(R\apart0\),
construct \(\phi=\atanTwo(B,A)\) with its quadrant certificate; then
\(\cos\phi=A/R\) and \(\sin\phi=B/R\).  If \(R=0\) is positively certified,
the left side is identically zero and no phase is needed.  What is not
constructively valid is an unguarded algorithm that decides which of these two
cases holds for an arbitrary real pair.

The tangent action of a rotation is the Mobius map
\[
 M_a(X)=\frac{\cos a\,X-\sin a}{\sin a\,X+\cos a}.
 \tag{11.4}
\]
Matrix multiplication and the addition formulas prove
\[
 M_a(M_b(X))=M_{a+b}(X),
 \tag{11.5}
\]
where all evaluated denominators are apart from zero.  Cross multiplication is
a polynomial identity and is one of the exact certificates.

\section{Euler, hyperbolic, series, derivatives, and integrals}
\label{sec:analytic-families}

Euler's formulas are definitions in the chosen construction:
\[
 \cos x=\frac{e^{ix}+e^{-ix}}2,
 \qquad
 \sin x=\frac{e^{ix}-e^{-ix}}{2i}.
 \tag{12.1}
\]
Defining
\(\cosh x=(e^x+e^{-x})/2\) and
\(\sinh x=(e^x-e^{-x})/2\) gives without further analysis
\[
 \cos(ix)=\cosh x,
 \quad \sin(ix)=i\sinh x,
 \quad \cosh(ix)=\cos x,
 \quad \sinh(ix)=i\sin x.
 \tag{12.2}
\]
The factorial definitions immediately yield the Maclaurin series
\begin{align}
 \sin x&=\sum_{n=0}^\infty(-1)^n\frac{x^{2n+1}}{(2n+1)!},\tag{12.3}\\
 \cos x&=\sum_{n=0}^\infty(-1)^n\frac{x^{2n}}{(2n)!},\tag{12.4}
\end{align}
with a modulus obtained from (3.2).  Series for tangent, secant, and related
functions are obtained by formal power-series inversion on a neighborhood with
a certified nonzero constant term; evaluation at a real point additionally
requires a denominator apartness bound.

\subsection{Constructive differentiation}

For bounded \(x\) and small \(h\), the exponential series gives the explicit
remainder estimate
\[
 \norm{e^h-1-h}_1\le C_B\norm h_1^2,
 \tag{12.5}
\]
where, for example, \(C_B=2^{2B+4}\) is a primitive-recursive integer
majorant on \(\norm h_1\le1\).  Using
\(e^{x+h}-e^x=e^x(e^h-1)\), equation (12.5) gives a primitive-recursive
modulus for
\[
 \frac{e^{x+h}-e^x}{h}\longrightarrow e^x.
 \tag{12.6}
\]
Taking real and imaginary parts proves
\[
 (\sin x)'=\cos x,
 \qquad (\cos x)'=-\sin x.
 \tag{12.7}
\]
Product and quotient rules follow from finite algebra plus a lower bound on the
quotient denominator.  Thus
\[
 (\tan x)'=\sec^2x,
 \quad (\cot x)'=-\csc^2x,
 \quad (\sec x)'=\sec x\tan x,
 \quad (\csc x)'=-\csc x\cot x
 \tag{12.8}
\]
on constructively open regions carrying the relevant apartness witnesses.
Termwise integration of (12.3)--(12.4), with the same factorial tail modulus,
gives
\[
 \int_a^b\sin x\,dx=\cos a-\cos b,
 \qquad
 \int_a^b\cos x\,dx=\sin b-\sin a
 \tag{12.9}
\]
for rational or rapid-name endpoints.  A general constructive Riemann-integration
library is not needed for these two identities; the integrated series itself is
a certified antiderivative.

\section{Principal inverse functions and branch-safe identities}
\label{sec:inverse}

Classical tables often write inverse identities without enough branch data.
Here the principal inverses are algorithms, and their uniqueness intervals are
part of every theorem.

\subsection{Constructing \(\arcsin\)}

On \([-q,q]\), sine is increasing.  A quantitative inverse modulus is available
even at the endpoints.  If \(-q\le a<b\le q\), put \(d=b-a\).  The
sum-to-product formula gives
\[
 \sin b-\sin a
 =2\cos\frac{a+b}{2}\sin\frac d2.
 \tag{13.1}
\]
For fixed \(d\), the smallest possible cosine factor occurs at an endpoint and
is at least \(\sin(d/2)\).  With (4.4),
\[
 \sin b-\sin a\ge2\sin^2\frac d2\ge\frac{d^2}{18}.
 \tag{13.2}
\]
Thus a sine-value uncertainty below \(2^{-2p-C}\) locates its argument to
within \(2^{-p}\).

For \(x\in[-1,1]\), a bounded rational bisection of \([-q,q]\) compares a
certified interval for \(\sin m-x\) with zero.  A separated sign keeps the
appropriate half; an interval containing zero invokes (13.2) and already gives
a sufficiently accurate argument.  The resulting rapid name is
\(\arcsin x\).  The construction is primitive recursive because the number of
steps and every requested precision are explicit functions of \(p\).

Define
\[
 \arccos x=q-\arcsin x,
 \qquad
 \arctan x=\arcsin\frac{x}{\sqrt{1+x^2}}.
 \tag{13.3}
\]
The square-root denominator is at least one, so it is uniformly apart from
zero.

\begin{theorem}[Principal compositions]
Constructively,
\begin{align*}
 \sin(\arcsin x)&=x &&(-1\le x\le1),\\
 \arcsin(\sin y)&=y &&(-q\le y\le q),\\
 \cos(\arccos x)&=x &&(-1\le x\le1),\\
 \arccos(\cos y)&=y &&(0\le y\le2q),\\
 \tan(\arctan x)&=x &&(x\in\R),\\
 \arctan(\tan y)&=y &&(-q<y<q\text{ with endpoint separation}).
\end{align*}
\end{theorem}

\begin{proof}
The forward equations are the specifications of the bisection algorithms.
The reverse equations use strict monotonicity and the stated interval
membership.  Tangent is handled through (13.3), the positive square root, and
the half-plane in which cosine is positive.  All endpoint inequalities are
finite interval certificates.
\end{proof}

\subsection{Addition of inverse tangents}

The globally correct branch-safe statement is
\[
 \arctan x+\arctan y
 =\atanTwo(x+y,1-xy),
 \tag{13.4}
\]
where the right side carries a certificate selecting the argument of the vector
\((1-xy,x+y)\), and the vector is apart from \((0,0)\).  If a branch certificate
shows \(1-xy>0\), this reduces to
\[
 \arctan x+\arctan y
 =\arctan\frac{x+y}{1-xy}.
 \tag{13.5}
\]
If it is negative, a certified \(\pm\pi\) correction is required according to
the sign of \(x+y\).  The proof applies the tangent addition formula and then
uses uniqueness in the selected open interval.  This also supplies a clean
constructive proof of Machin-type formulas: multiply the corresponding Gaussian
rationals, verify the resulting argument exactly, and certify the quadrant.

\section{Triangle identities}
\label{sec:triangles}

Let \(\alpha+\beta+\gamma=\pi\).  Substitute
\(\gamma=\pi-\alpha-\beta\) and use the addition and reflection formulas.
Every polynomial triangle identity is then a Laurent certificate in two angle
variables.  Representative examples are
\begin{align}
 \tan\alpha+\tan\beta+\tan\gamma
 &=\tan\alpha\tan\beta\tan\gamma,\tag{14.1}\\
 \cot\beta\cot\gamma+\cot\gamma\cot\alpha+
 \cot\alpha\cot\beta&=1,\tag{14.2}\\
 \sin\alpha+\sin\beta+\sin\gamma
 &=4\cos\frac\alpha2\cos\frac\beta2\cos\frac\gamma2,\tag{14.3}\\
 \cos\alpha+\cos\beta+\cos\gamma
 &=1+4\sin\frac\alpha2\sin\frac\beta2\sin\frac\gamma2,\tag{14.4}\\
 \sin^2\alpha+\sin^2\beta+\sin^2\gamma
 &=2+2\cos\alpha\cos\beta\cos\gamma,\tag{14.5}\\
 \cos^2\alpha+\cos^2\beta+\cos^2\gamma
 &=1-2\cos\alpha\cos\beta\cos\gamma.\tag{14.6}
\end{align}
Quotient formulas (14.1)--(14.2) require the relevant sine and cosine
apartness witnesses.  Equations (14.3)--(14.6) are unconditional.  The exact
ledger contains seventeen triangle and half-angle variants, including signed
and double-angle forms.

As an illustration, (14.3) follows by first combining \(\sin\alpha+\sin\beta\):
\[
 2\sin\frac{\alpha+\beta}{2}\cos\frac{\alpha-\beta}{2}
 =2\cos\frac\gamma2\cos\frac{\alpha-\beta}{2}.
\]
Adding \(\sin\gamma=2\sin(\gamma/2)\cos(\gamma/2)\), factoring
\(2\cos(\gamma/2)\), and applying the cosine sum formula gives the right side.
This proof uses no positivity assumption; positivity is needed only when a
square root or quotient version is selected.

\section{Infinite dyadic products and Viete's formula}
\label{sec:infinite}

Put \(t=x/2^N\) in (9.11).  The finite identity becomes
\[
 \prod_{j=1}^N\cos\frac{x}{2^j}
 =\frac{\sin x}{2^N\sin(x/2^N)},
 \tag{15.1}
\]
when the denominator is apart from zero, with an unconditional
cross-multiplied version.

For \(\abs t\le1\), the alternating series gives
\[
 \abs{\sin t-t}\le\frac{\abs t^3}{6}.
 \tag{15.2}
\]
Therefore
\[
 \abs{2^N\sin(x/2^N)-x}
 \le\frac{\abs x^3}{6\,4^N}.
 \tag{15.3}
\]
Given a magnitude bound for \(x\) and a desired precision \(p\), a bounded
search for \(N\) with \(4^N>2^{p+C}\abs x^3\) is primitive recursive.  Hence
(15.1) has the constructive limit
\[
 \frac{\sin x}{x}=\prod_{j=1}^\infty\cos\frac{x}{2^j},
 \tag{15.4}
\]
with quotient notation guarded by \(x\apart0\); at \(x=0\) the continuous value
is one and can be stated directly.

Taking \(x=q=\pi/2\) gives
\[
 \frac{2}{\pi}
 =\prod_{j=1}^\infty\cos\frac{\pi}{2^{j+1}}.
 \tag{15.5}
\]
Let \(r_0=0\) and \(r_{n+1}=\sqrt{2+r_n}\), always choosing the nonnegative
square root.  Repeated positive half-angle formulas prove
\(\cos(\pi/2^{n+2})=r_{n+1}/2\), so
\[
 \frac2\pi=\prod_{n=1}^\infty\frac{r_n}{2},
 \tag{15.6}
\]
which is the nested-radical Viete formula.  Equations (15.3) and the square-root
modulus make (15.6) a theorem of the interpreted \(\iPRA\), not merely a
classical limit assertion.

The full Euler--Weierstrass factorization
\(\sin x=x\prod_{n\ge1}(1-x^2/(n^2\pi^2))\) is intentionally not claimed here.
Its constructive proof needs a separate theory of normally convergent infinite
products and a zero-classification argument.  None of the finite or dyadic
identities above depends on it.

\section{Bessel generating identities: a constructive extension path}
\label{sec:bessel}

Several entries in the source list express compositions through Bessel
functions.  They are not part of the 227 finite certificates, but their standard
proof is compatible with \(\iPRA\).  For \(n\ge0\), define
\[
 J_n(z)=\sum_{m=0}^\infty
 \frac{(-1)^m}{m!(m+n)!}\left(\frac z2\right)^{2m+n},
 \qquad J_{-n}(z)=(-1)^nJ_n(z).
 \tag{16.1}
\]
Factorial ratios again give an explicit geometric tail after a
primitive-recursive index.  Expanding two exponentials and collecting the
coefficient of a formal Laurent variable \(t\) yields
\[
 \exp\left(\frac z2(t-t^{-1})\right)
 =\sum_{n\in\Z}J_n(z)t^n.
 \tag{16.2}
\]
For every requested precision, only a finite rectangle in the double series is
needed; the omitted rows and columns have factorial tail bounds of the form in
Lemma \ref{lem:exp-tail}.  Thus rearrangement is justified by explicit finite
error estimates rather than a classical unconditional-convergence theorem.
Setting \(t=e^{i\theta}\) gives
\[
 e^{iz\sin\theta}=\sum_{n\in\Z}J_n(z)e^{in\theta},
 \tag{16.3}
\]
and replacing \(t\) by \(ie^{i\theta}\) gives the cosine form
\[
 e^{iz\cos\theta}=\sum_{n\in\Z}i^nJ_n(z)e^{in\theta}.
 \tag{16.4}
\]
A complete Lean implementation needs a reusable two-dimensional series-tail
library, but no stronger logical principle.  The formalization plan in
Section \ref{sec:lean} lists this after the finite and one-dimensional analytic
layers.

\section{Why the proofs are intuitionistic and primitive recursive}
\label{sec:logical-audit}

The following audit records where classical-looking steps were replaced.

\begin{longtable}{@{}p{0.25\textwidth}p{0.68\textwidth}@{}}
\toprule
Classical shorthand & Constructive replacement\\
\midrule
\endfirsthead
\toprule
Classical shorthand & Constructive replacement\\
\midrule
\endhead
``Take a real limit'' & A rapid dyadic name with a displayed primitive-recursive
modulus and a proof of every precision inequality.\\
``The denominator is nonzero'' & A finite apartness witness giving a positive
dyadic lower bound.\\
``Choose the sign of a square root'' & A quadrant or interval-location
certificate; zero endpoints are separate exact cases.\\
Intermediate value theorem & A bounded located sign-search with quantitative
monotonicity estimates.\\
Monotonic inverse & Bisection plus an explicit inverse modulus such as (13.2).\\
Rearrange an infinite series & Compare finite rectangles and bound every omitted
tail by a primitive-recursive majorant.\\
Complex roots of unity & A finite polynomial quotient ring; no algebraic closure
or choice of arbitrary roots.\\
``For all computable reals'' & A uniform proof schema over fixed
primitive-recursive name terms and their arithmetic certificates.\\
Ring simplification & A finite Gaussian Laurent normalizer whose soundness is
proved by quantifier-free induction.\\
\bottomrule
\end{longtable}

All unbounded statements in the final formulas are ordinary first-order
universal quantifiers over natural precision parameters.  Induction is applied
to quantifier-free invariants for finite sums, products, recursion equations,
normalization lists, and bounded search.  The logic never uses excluded middle
for real equality, Markov's principle, countable choice, or a nonconstructive
completion of the rationals.

\section{Lean formalization plan}
\label{sec:lean}

The archive contains three Lean files.  They are deliberately separated by
proof layer.

\subsection{Layer A: abstract algebraic trigonometry}

\code{Lean/TrigAlgebra.lean} introduces a class
\begin{lstlisting}
class TrigSystem (Angle R) [AddCommGroup Angle] [CommRing R] where
  sin cos : Angle -> R
  q : Angle
  sin_zero : sin 0 = 0
  cos_zero : cos 0 = 1
  sin_add : sin (x+y) = sin x*cos y + cos x*sin y
  cos_add : cos (x+y) = cos x*cos y - sin x*sin y
  pythagorean : sin x ^ 2 + cos x ^ 2 = 1
  sin_q : sin q = 1
  cos_q : cos q = 0
\end{lstlisting}
It then derives parity, subtraction, quarter-turn shifts, periods, double and
triple angles, Chebyshev recurrences, product-to-sum identities, and half-angle
squares by \code{ring} or \code{linear_combination}.  This layer should be made
fully compiling first.  It contains no constructive-real implementation and no
classical analysis.

\subsection{Layer B: exact certificate syntax}

\code{Lean/LaurentCertificate.lean} fixes the intended types:
Gaussian rationals, finitely supported integer monomials, Laurent polynomials,
integer angle vectors, and numerator/denominator pairs.  Completion consists of
implementing the Python reference normalizer and proving:
\begin{enumerate}
\item normalization preserves evaluation;
\item normalized zero evaluates to zero;
\item \code{cis}, \code{sine}, and \code{cosine} syntax evaluates correctly in
any \code{TrigSystem};
\item a rational certificate is sound under denominator apartness.
\end{enumerate}
All inductions are structural or on finite maps.  No tactic beyond finite ring
normalization is logically needed.

\subsection{Layer C: direct integration with ConciseIPRA}

\code{Lean/IPRABridge.lean} imports
\code{ConciseIPRA.PR} and \code{ConciseIPRA.EqArith}.  The implementation order
is:
\begin{enumerate}
\item Define signed integers and dyadics as explicit \code{PR} symbols; prove
all computation equations through \code{PRA.proj}, \code{PRA.comp},
\code{PRA.precZero}, and \code{PRA.precSucc}.
\item Define formulas \code{Regular}, \code{EqReal}, \code{Bounds}, and
\code{Apart} over \code{PTm}; implement name constructors for negation,
addition, multiplication, reciprocal, and square root.
\item Prove regularity closure with \code{PRA.qfInd}.  Keep the magnitude bound
as part of every name package so multiplication and series precision shifts are
primitive-recursive terms.
\item Define finite complex exponential sums and the index
\code{expTerms B p = 2*B*B + 2*B + p + 4}.  Formalize Lemma
\ref{lem:exp-tail}.
\item Prove the finite coefficient identity (3.6), then exponential addition,
sine/cosine addition, and Pythagoras.
\item Implement the bounded sign-search of Section \ref{sec:pi}; prove the
quarter-turn specification.
\item Connect the Laurent evaluator to the real-name operations and import the
227 certificates as data, or regenerate them inside Lean.
\item Add inverse bisection, derivative moduli, and the dyadic infinite product.
\item Finally add the two-dimensional factorial-tail library for Bessel
identities.
\end{enumerate}

The existing ProveIt development already supplies the crucial syntactic objects:
\code{PR k}, \code{PTm n}, the axiom family \code{PRAAx}, and the theorem
\code{PRA.qfInd}.  The bridge should reuse these rather than formalizing a second
arithmetic.

\subsection{Suggested theorem interfaces}

The following interfaces make the path to completion unambiguous.
\begin{lstlisting}
-- arithmetic names
regular_add  : Regular a -> Regular b -> Regular (add a b)
regular_mul  : Bounds a A -> Bounds b B ->
               Regular a -> Regular b -> Regular (mul a b)
regular_inv  : Apart a w -> Regular a -> Regular (inv a w)

-- analytic core
exp_tail     : TailBound B p (expTerms B p)
cexp_add     : EqComplex (cexp (addC x y)) (mulC (cexp x) (cexp y))
pythagorean  : EqReal (sin x * sin x + cos x * cos x) one
quarter_spec : EqReal (cos quarterTurn) zero
quarter_sin  : EqReal (sin quarterTurn) one

-- certificate bridge
normalize_sound : eval (normalize P) rho = eval P rho
certificate_sound : check lhs rhs = true -> Guards rho lhs rhs ->
                    EqReal (evalRF lhs rho) (evalRF rhs rho)
\end{lstlisting}
The actual types will use arithmetized derivations
\(\Gamma\vdash[\mathsf{PRAAx}]A\), not meta-level Lean propositions, for the
final soundness theorem.

\section{Executable verification and reproducibility}
\label{sec:verification}

Run
\begin{lstlisting}
./scripts/verify_all.sh
\end{lstlisting}
from the archive root.  The script constructs the Gaussian Laurent model,
normalizes both polynomial and cross-multiplied rational residues, and writes
\code{artifacts/identity_certificates.json}.  It uses Python's exact
\code{Fraction}; no numerical sampling is accepted as a proof.

\VerbatimInput[fontsize=\small]{artifacts/verification_report.txt}

The checked instances include: basic and Pythagorean identities; parity,
reflection, shift, and periodicity; all addition and common quotient formulas;
double, triple, and binomial multiple-angle formulas; Chebyshev recurrences;
half-angle and Weierstrass formulas; product/sum conversions; general products
through six factors; power reductions through fifth powers; elementary-
symmetric many-angle formulas through four variables; Lagrange and Dirichlet
families through ten; arithmetic progressions through eight; finite Viete
products through nine; a linear-fractional rotation; amplitude squares;
Hermite identities for two and three poles; seventeen triangle identities; and
Euler's formulas.

The exact checker has a narrower role than the report's proof.  It detects
transcription and sign errors in a large finite catalog.  The arbitrary-\(n\)
claims are justified by the finite induction or binomial arguments in the text,
not by extrapolation from checked instances.

\section{Limitations and further research questions}

\begin{enumerate}
\item Complete the direct arithmetization of rapid names on top of
\code{ConciseIPRA.PR}, replacing every \code{axiom} in
\code{Lean/IPRABridge.lean} with a derivation object.
\item Prove the Laurent normalizer correct in Lean and import the JSON ledger as
certificates.  A useful comparison is whether direct reflection or proof-term
generation gives the smaller trusted kernel footprint.
\item Develop a constructive normally-convergent product library and use it to
formalize the Euler--Weierstrass sine product in \(\iPRA\) with an explicit
uniform modulus away from its zeros.
\item Add an efficient interval \(\atanTwo\) whose branch certificate composes
well under Gaussian multiplication, then generate a large library of
Machin-type arctangent identities automatically.
\item Formalize the two-dimensional factorial-tail lemma and Jacobi--Anger
identities.  Determine the smallest reusable series interface sufficient for
both Bessel and trigonometric applications.
\item Extract a standalone proof checker in the style of ConciseIPRA's scalar
calculi.  The algebraic checker is already tiny; the research question is how
small a checker can validate the analytic modulus certificates as well.
\item Compare the exponential-series construction of \(\pi\) here with
polygonal, arctangent, and AGM constructions under the metric ``size of a
straightforward iPRA proof checker.''
\end{enumerate}

\appendix
\section{A compact catalog of proved formulas}

This appendix groups formulas by their proof generator.  All quotient signs
mean ``under apartness of every displayed denominator.''

\subsection*{Generated from \(\cis(x+y)=\cis x\cis y\)}
\begin{align*}
\sin(x\pm y)&=\sin x\cos y\pm\cos x\sin y,\\
\cos(x\pm y)&=\cos x\cos y\mp\sin x\sin y,\\
\tan(x\pm y)&=\frac{\tan x\pm\tan y}{1\mp\tan x\tan y},\\
\cot(x+y)&=\frac{\cot x\cot y-1}{\cot x+\cot y},\\
\cot(x-y)&=\frac{\cot x\cot y+1}{\cot y-\cot x},\\
\sec(x+y)&=\frac{\sec x\sec y}{1-\tan x\tan y},\\
\csc(x+y)&=\frac{\sec x\sec y}{\tan x+\tan y}.
\end{align*}

\subsection*{Generated by binomial expansion}
Equations (7.1)--(7.2), all double/triple-angle polynomial forms, all guarded
quotient multiple-angle forms, Chebyshev formulas (7.4), and power reductions
(8.6)--(8.7).

\subsection*{Generated by the two-factor product identities}
Equations (9.1)--(9.6), all sum-to-product formulas, triangle formulas
(14.3)--(14.6), and the signed variants in the certificate ledger.

\subsection*{Generated by finite geometric sums}
Equations (10.2)--(10.8), arbitrary arithmetic-progression sums, the Dirichlet
kernel, odd-harmonic cosine sums, and their sine/cosine shifted variants.

\subsection*{Generated by repeated doubling}
Finite Viete products (9.11), the dyadic product (15.4), Viete's nested radical
(15.6), and their cross-multiplied zero-safe versions.

\section{Artifact inventory}

\begin{description}
\item[\code{trigonometric_identities_ipra.tex}] This source.
\item[\code{trigonometric_identities_ipra.pdf}] Compiled report.
\item[\code{scripts/verify_identities.py}] Exact Gaussian Laurent checker.
\item[\code{scripts/verify_all.sh}] Reproducibility wrapper.
\item[\code{artifacts/identity_certificates.json}] Per-identity zero-residue
ledger.
\item[\code{artifacts/verification_report.txt}] Summary generated by the
checker.
\item[\code{artifacts/coverage.csv}] Coverage and constructive side-condition
matrix.
\item[\code{Lean/TrigAlgebra.lean}] Abstract algebraic Lean prototype.
\item[\code{Lean/LaurentCertificate.lean}] Certificate datatype and soundness
outline.
\item[\code{Lean/IPRABridge.lean}] Direct ConciseIPRA implementation blueprint.
\end{description}

\begin{thebibliography}{9}
\bibitem{wiki}
Wikipedia contributors, \emph{List of trigonometric identities},
\url{https://en.wikipedia.org/wiki/List_of_trigonometric_identities}, accessed
15 November 2026.

\bibitem{proveit}
V. Reshetnikov, \emph{ProveIt}, especially
\url{https://github.com/VladimirReshetnikov/ProveIt/tree/main/Logic/PrimitiveRecursiveArithmetic/ConciseIPRA}
and
\url{https://github.com/VladimirReshetnikov/ProveIt/tree/main/Analysis/TrigonometricIdentities}.

\bibitem{bishop}
E. Bishop and D. Bridges, \emph{Constructive Analysis}, Springer, 1985.

\bibitem{bridges-richman}
D. Bridges and F. Richman, \emph{Varieties of Constructive Mathematics},
Cambridge University Press, 1987.
\end{thebibliography}

\end{document}
